% Carried verbatim from the current consensus abstract in OpenGU-DocMap.
% Its findings sentence remains provisional pending scientific review.
\begin{abstract}
Graph unlearning (GU) methods are being investigated as a family of machine unlearning techniques for protecting data privacy. These methods face distinct challenges due to interdependencies among nodes and the underlying graph topology. However, red-teaming of GU methods through strategically selected deletion requests remains underexplored. More specifically, the different components of GU performance changes under such attacks have received limited analysis. To address these gaps, we propose a framework from an adversarial perspective. It uses deletion-target selectors to construct adversarial requests and examines how GU methods respond under these attacks. Within this framework, we define attacker access levels, including gray-box and white-box settings, based on the information available for deletion-target selection. Our extensive experiments demonstrate how performance decomposition and complementary metrics distinguish GU responses to deletion requests. Our case studies reveal method- and dataset-dependent responses, with substantial utility degradation occurring even under random deletion and targeted requests producing divergent effects on GU and full retraining. The framework facilitates selector design and informs the development of more robust GU methods.
\end{abstract}
